\section{Discussion}
\label{sec:discussion}

We state the design knowledge as four principles, each with its evidence and boundary condition. For each we say which evidence is registered; results ``added after review'' are outside the pre-registered plan.

\paragraph{Principle 1: execute in code the rules that stop or hold the checklist}
On the \sthreeClaimsDenied{} denied-event and \sthreeClaimsOpenEvent{} open-event tenancy claims, \casepath made \sthreeCasePathDeniedStrict{} and \sthreeCasePathOpenEventStrict{} strict unsafe errors per claim, Rules in context \sthreeRicDeniedStrict{} and \sthreeRicOpenEventStrict, and the frontier single prompt \sthreeRicFrontierDeniedStrict{} and \sthreeRicFrontierOpenEventStrict{} (registered subset rows, Table~\ref{tab:s3-subsets}). On the same claims Model executes rules, which states the semantics in its prompt but runs no code, made \sthreeMerDeniedStrict{} and \sthreeMerOpenEventStrict: stating the semantics closed most of the gap, and executing them in code closed most of the rest. \emph{Boundary:} the gain scales with the share of such claims, which the test set chose; on ordinary claims the product did not differ significantly from the frontier model.

\paragraph{Principle 2: give an open fact its own value, and evaluate every completion}
Registered: the product held \sthreeCasePathTenCorrectHold\% of the documents the tenancy references hold, against \sthreeRicTenCorrectHold\% and \sthreeRicFrontierTenCorrectHold\% for the single prompts; two-valued execution of its own readings made \sthreeTwoValFalseTenStrict{} (unresolved taken as false) and \sthreeTwoValTrueTenStrict{} (taken as true) strict unsafe errors per claim, against \sthreeCasePathTenStrict. Added after review: against strong Kleene evaluation of the same facts, exact deferral made no fewer strict unsafe errors (\sthreeRElevenBExactStrict{} against \sthreeRElevenBKleeneStrict) but fewer deviations (\sthreeRElevenBExactAllDev{} against \sthreeRElevenBKleeneAllDev{} per claim), because strong Kleene evaluation held more documents that are required now or not needed. \emph{Boundary:} exact deferral never ran on theft, and there ``unresolved = false'' gives the product's requests, so the principle is supported on tenancy only, within the scope of Section~\ref{sec:rule}. It costs \sthreeCasePathTenSpurHold{} unneeded holds per claim, against \sthreeRicTenSpurHold{} and \sthreeRicFrontierTenSpurHold, and \sthreeCasePathDeniedSpurHold{} on denied-event claims, where the defining event stayed unresolved.

\paragraph{Principle 3: admit only agreeing readings with quotations found in the file}
No registered test isolates this principle. Added after review: of \sthreeRThreeCasePathFactRuns{} tenancy fact values the product used, \sthreeRThreeCasePathAdmittedWrong{} admitted values and \sthreeRThreeCasePathPresumedWrong{} presumptions were wrong. In development, a second reading that must agree cut theft errors (Table~\ref{tab:history}, recomputed by pairing readings across runs). \emph{Boundary:} in every wrong admission, all three readings agreed on the wrong value. The readings share one model and prompt, so agreement misses a shared misreading; whether a document is present is a model judgment outside this rule.

\paragraph{Principle 4: run the agents that cannot decide the checklist only where handlers need them}
Registered: no need added by an agent entered a scored checklist, and dropping the brief, the process builder and the result checker showed no detectable difference in errors at \sthreeCostChecklistShareTen\% (tenancy) and \sthreeCostChecklistShareTh\% (theft) of the full product's cost (Section~\ref{sec:test-results}). Added after review: the two settings' checklists differed no more than two runs of the same setting (Section~\ref{sec:test-review}). \emph{Boundary:} the two settings shared cached responses, and the brief and reply draft may be worth their cost and time to a handler, which we did not measure.

\paragraph{Where it does not pay}
On theft, the product's difference from Self-consistency did not reach significance. No hypothesis concerned \casepath, checklist only, on the frontier model (one run), so the following is descriptive. On theft it was not better than the frontier single prompt (\sthreeSecSevenAA{} against \sthreeSecSevenAB, $p\sthreeSecSevenAPeq$): the difference from the same model's single prompt was \sthreeHOneDiff{} errors per pair with the cost-efficient model and \sthreeSecSevenADiff{} with the frontier model. On tenancy the difference was similar for both models (\sthreeHTwoDiff{} and \sthreeSecSevenBDiff{} strict unsafe errors per claim against each model's single prompt). With the original rules the frontier single prompt made \sthreeRicOrigFrontierThErrors{} theft errors per pair, close to \casepath's \sthreeCasePathThErrors. We conjecture that execution pays where many rules are in force; the ShARC pilot (Section~\ref{sec:history}) is consistent with this but does not test it.

\paragraph{For an organisation}
An organisation must write workflows and presumptions and have them reviewed. Here two AI editor sessions reviewed the \rvRules{} theft rules; they changed 9 and queued 14 for a jurist, and no jurist has reviewed them; the review's duration was not recorded. Every presumption must cite a clause of the organisation's own sources (Section~\ref{sec:format}), so the owner of that clause answers for the presumption; in our tests the researchers wrote the clause. A wrong presumption drops the documents that depend on its fact, and the checklist names the clause, so the error can be traced. An auditor gets each deciding value with its quotation, each presumption with its clause, each hold with its question, and the hash of the code that ran. Holds add work: a held document waits for the claimant's answer to a question.

\paragraph{Limits}
The texts have 77--320 words, so long-file paths were never exercised. The tenancy references come from the engine \casepath implements, under a silence clause the researchers wrote, not the insurer: they test how a design applies the rules, not whether rules or presumptions are right, and that an announced but unreceived notice needs no document is a choice of the data design. Wrong presumptions were tested only by a stress test added after review: with 5, 10 and 20\% of unmentioned facts true, the product's errors per claim rose to \sthreeRFourFiveCasePath, \sthreeRFourTenCasePath{} and \sthreeRFourTwentyCasePath, and its difference from the frontier single prompt stayed at \sthreeRFourTwentyDiffRicFrontier. Theft reuses Study~1's conditions, with rules reviewed knowing them, so it tests new texts, not new rules; with \sthreeTheftClusters{} clusters, H4's failure does not show the designs equal. H3 was tested only because the order of the hypotheses changed shortly before the freeze. H1 to H3 hold for the cost-efficient model only.

\paragraph{The expert study}
A study with claims experts has not run. Its protocol, version 1.0-rc, will be frozen before the first session (Appendix~\ref{app:protocol}). Five of the insurer's claims experts will work through three closed files, each cut before the insurer's first request for documents. It asks how far the product agrees with that request (P), and what experts change and why. P is transcribed and mapped to the catalogue by the product builder's automated labeller before any run on the cut copy, and checked by the study lead, an author who is not a claims expert. Its limits: an interface unlike the tested commit's, cases seen in development, experts who may have handled a file, and one test whose smallest $p$ is $1/32$.
